\chapter{REVISÃO SISTEMÁTICA DE LITERATURA}

Este capítulo apresenta o protocolo e os resultados da revisão de literatura (RSL) sobre HTTP/3, QUIC, desempenho e implementação em APIs REST e serviços web. A revisão busca identificar e sintetizar as evidências disponíveis para apoiar as questões de pesquisa definidas no Capítulo 1.

\section{Protocolo de pesquisa}

Esta revisão de literatura foi conduzida com base em diretrizes para revisões sistemáticas em Engenharia de Software descritas por Kitchenham e Charters (2007). O objetivo é identificar, selecionar e sintetizar estudos sobre HTTP/3, QUIC, desempenho e implementação em APIs ou serviços web que ajudem a responder às questões de pesquisa.

A revisão seguiu as etapas de: (i) definição da questão de revisão e dos critérios de busca; (ii) execução da busca nas bases selecionadas; (iii) remoção de duplicatas; (iv) triagem por título e resumo; (v) leitura integral dos estudos potencialmente elegíveis; (vi) extração de dados e síntese temática dos estudos incluídos.

A questão norteadora da revisão foi ``\textit{Quais evidências a literatura apresenta sobre desempenho, implementação e migração de HTTP/3 e QUIC em APIs REST e serviços web, e quais lacunas permanecem quanto à adoção em ambientes Node.js?}\textit{''}

A \textit{string} de busca foi construída a partir de três blocos temáticos, combinados com o operador lógico AND:

\begin{enumerate}[label=\alph*)]
\item protocolo: ("HTTP/3"OR "HTTP 3"OR HTTP3 OR QUIC);
\item contexto de API: ("REST"OR "RESTful"OR "web API"OR "Application Programming Interface");
\item aspecto investigado: (performance OR latency OR migration OR implementation OR "Node.js"OR "Node js").
\end{enumerate}

A string completa utilizada como referência foi:

("HTTP/3"OR "HTTP 3"OR HTTP3 OR QUIC) AND ("REST"OR "RESTful"OR"web API"OR "Application Programming Interface") AND (performance OR latency OR migration OR implementation OR "Node.js"OR "Node js")

Foram aplicados filtros de publicação entre 2018 e 2026, período que abrange a consolidação do QUIC e do HTTP/3 como padrões IETF. A Tabela 2 apresenta a adaptação da \textit{string} conforme a sintaxe exigida por cada base de dados. As buscas foram executadas em 29 de junho de 2026.

\begin{table}[htbp]
\centering
\caption{Adaptação da string de busca por base de dados}
\small
\setlength{\tabcolsep}{3pt}
\begin{tabular}{@{}>{\raggedright\arraybackslash}p{32mm}>{\raggedright\arraybackslash}p{112mm}@{}}
\hline
\textbf{Base} & \textbf{String adaptada (2018--2026)} \\
\hline
IEEE Xplore & ("HTTP/3"OR "HTTP 3"OR HTTP3 OR QUIC) AND ("REST"OR "RESTful"OR "web API") AND (performance OR latency OR migration OR implementation OR "Node.js") \\
ACM Digital Library & AllField: ("HTTP/3"OR QUIC) AND (REST OR "web API") AND (performance OR latency OR migration OR implementation) \\
Scopus & TITLE-ABS-KEY(|("HTTP/3"OR HTTP3 OR QUIC) AND ("REST"OR "web API") AND (performance OR latency OR migration OR implementation)) AND PUBYEAR > 2017 AND PUBYEAR < 2027 \\
Springer Link & ("HTTP/3"OR "HTTP 3"OR HTTP3 OR QUIC) AND (REST OR RESTful OR "web API") AND (performance OR latency OR migration OR implementation) \\
Google Scholar & "HTTP/3"OR QUIC REST API performance latency migration Node.js (2018..2026) \\
\hline
\end{tabular}
\fonte{Fonte: elaborado pelo autor}
\end{table}

Foram consultadas cinco bases de dados, selecionadas por cobrirem periódicos e conferências das áreas de redes, sistemas distribuídos e Engenharia de Software:

\begin{enumerate}[label=\alph*)]
\item IEEE Xplore - conferências e periódicos IEEE em redes e sistemas (ex.: CoNTESA, ICMI);
\item ACM Digital Library - publicações ACM em sistemas distribuídos, web e protocolos;
\item Scopus - indexação ampla de periódicos e conferências internacionais;
\item Springer Link - artigos em ciência da computação e redes;
\item Google Scholar - busca complementar, com deduplicação rigorosa em relação às demais bases.
\end{enumerate}

Os critérios de seleção foram definidos antes da triagem dos estudos, conforme descrito a seguir.

Critérios de inclusão:
\begin{enumerate}[label=\alph*)]
\item artigos completos, extended abstracts de conferências ou normas técnicas (RFCs) relevan-tes, com revisão por pares ou publicação em organismos de padronização;
\item estudos que abordem HTTP/3, QUIC, desempenho, latência, migração ou implementação em APIs, serviços web ou arquiteturas de backend;
\item publicações em inglês ou português;
\item período de publicação entre 2018 e 2026;
\item acesso ao texto completo ou resumo suficiente para avaliação de elegibilidade.
\end{enumerate}

Critérios de exclusão:
\begin{enumerate}[label=\alph*)]
\item duplicatas identificadas entre bases (mesmo DOI, título ou autores);
\item editoriais, keynotes ou trabalhos sem contribuição técnica verificável;
\item estudos focados exclusivamente em frontend, streaming de mídia ou domínios sem relação com APIs REST/backend;
\item pré-prints sem revisão por pares, exceto quando citados como complemento a versões publicadas.
\end{enumerate}

O processo de seleção dos estudos seguiu um funil em cinco etapas:

\begin{enumerate}[label=\alph*)]
\item identificação - registros retornados pelas buscas em cada base (Tabela 4);
\item deduplicação - remoção de 515 registros duplicados, principalmente entre Scopus, Google Scholar e Springer Link, resultando em 284 estudos únicos;
\item triagem - leitura de título e resumo, com exclusão de 260 estudos por não atenderem aos critérios de inclusão (foco em domínios distintos, ausência de relação com HTTP/3/QUIC ou APIs);
\item elegibilidade - leitura integral de 24 estudos, dos quais 13 foram excluídos por não apresentarem dados comparativos, implementação ou análise aplicável ao escopo deste TCC;
\item inclusão - 11 documentos selecionados para síntese: duas normas técnicas (RFC 9000 e RFC 9114) e nove estudos empíricos ou de revisão com dados comparativos (Tabelas 5 e 6).
\end{enumerate}

A Tabela 3 resume o quantitativo de estudos em cada etapa do processo de seleção.


\begin{table}[htbp]
\centering
\caption{Quantitativo de estudos em cada etapa da revisão}
\small
\setlength{\tabcolsep}{3pt}
\begin{tabular}{@{}>{\raggedright\arraybackslash}p{100mm}>{\raggedright\arraybackslash}p{28mm}@{}} \hline
\textbf{Etapa} & \textbf{Registros} \\
\hline
Identificação (total bruto) & 799 \\
Após deduplicação & 284 \\
Após triagem (título e resumo) & 24 \\
Após elegibilidade (leitura integral) & 11 \\
\hline
\end{tabular}
\fonte{Fonte: elaborado pelo autor (buscas realizadas em 29/06/2026). Nota: os 11 documentos incluem 2 normas (RFC) e 9 estudos empíricos ou de revisão sintetizados nos eixos temáticos.}
\end{table}

A Tabela 4 apresenta a quantidade de registros identificados em cada base de dados consultada, antes da deduplicação. O total bruto de 799 registros reflete sobreposição entre indexadores; após remoção de duplicatas, permaneceram 284 estudos únicos para triagem.


\begin{table}[htbp]
\centering
\caption{Quantidade de registros identificados por base de dados}
\small
\setlength{\tabcolsep}{3pt}
\begin{tabular}{@{}>{\raggedright\arraybackslash}p{88mm}>{\raggedright\arraybackslash}p{32mm}>{\raggedright\arraybackslash}p{24mm}@{}} \hline
\textbf{Base de dados} & \textbf{Registros encontrados} & \textbf{Estudos incluídos} \\
\hline
IEEE Xplore & 41 & 4 \\
ACM Digital Library & 27 & 2 \\
Scopus & 93 & 3 \\
Springer Link & 470 & 1 \\
Google Scholar & 168 & 1 \\
\textbf{Total (antes da deduplicação)} & \textbf{799} & --- \\
\textbf{Estudos únicos (após deduplicação)} & \textbf{284} & --- \\
\textbf{Documentos na síntese (2 RFC + 9 empíricos)} & --- & \textbf{11} \\
\hline
\end{tabular}
\fonte{Fonte: elaborado pelo autor (buscas realizadas em 29/06/2026). Nota: a coluna ``Estudos incluídos'' indica em qual base cada um dos 11 documentos finais foi originalmente identificado.}
\end{table}

Os quantitativos por base foram obtidos pelas interfaces de busca avançada de cada plataforma em 29/06/2026, com as strings da Tabela 2. As normas RFC 9000 e RFC 9114 não integram esses totais; foram incluídas manualmente por fundamentarem semanticamente QUIC e HTTP/3 (Tabela 5).

\section{Resultados da Revisão da Literatura}

Esta seção apresenta os resultados da revisão, organizados em base normativa e síntese empírica. As Tabelas 5 e 6 resumem os documentos incluídos após a elegibilidade.

\begin{table}[htbp]
\centering
\caption{Base normativa incluída na revisão (não contabilizada como evidência empírica)}
\small
\setlength{\tabcolsep}{3pt}
\begin{tabular}{@{}>{\raggedright\arraybackslash}p{52mm}>{\raggedright\arraybackslash}p{14mm}>{\raggedright\arraybackslash}p{78mm}@{}}
\hline
\textbf{Documento} & \textbf{Ano} & \textbf{Contribuição} \\
\hline
(QUIC Working Group -- IETF, 2021) & 2021 & RFC 9000 --- especificação QUIC (UDP, TLS 1.3, \textit{streams}) \\
(IETF, 2022) & 2022 & RFC 9114 --- mapeamento HTTP sobre QUIC \\
\hline
\end{tabular}
\fonte{Fonte: elaborado pelo autor.}
\end{table}

\begin{table}[htbp]
\centering
\caption{Estudos empíricos e de revisão incluídos na síntese temática (n = 9)}
\small
\setlength{\tabcolsep}{3pt}
\begin{tabular}{@{}>{\raggedright\arraybackslash}p{48mm}>{\raggedright\arraybackslash}p{12mm}>{\raggedright\arraybackslash}p{18mm}>{\raggedright\arraybackslash}p{66mm}@{}} \hline
\textbf{Estudo} & \textbf{Ano} & \textbf{Tipo} & \textbf{Contribuição principal} \\
\hline
(PERNA et al., 2022) & 2022 & Medição & Adoção e desempenho do HTTP/3 em larga escala na web \\
(TRIPATHI; MIRAZ; JOSHI, 2023) & 2023 & Medição & Latência HTTP/3 vs HTTP/1.1 em APIs RESTful \\
(CHATZOGLOU et al., 2023) & 2023 & Análise & Segurança do HTTP/3 comparado ao HTTP/2 \\
(KOCH; FALOWO; ELROD, 2024) & 2024 & Revisão & Estado da implementação e adoção do HTTP/3 \\
(ZHANG et al., 2025) & 2025 & Revisão & Desempenho, limitações e melhorias do HTTP/3 \\
(YU; BENSON, 2021) & 2021 & Medição & QUIC/HTTP/3 em produção (grandes provedores) \\
(TREVISAN et al., 2021) & 2021 & Medição & Adoção e desempenho do HTTP/3 em escala \\
(MUÑOZ et al., 2025) & 2025 & Medição & HTTP/3 vs HTTP/2 com \textit{proxy} \\
(ZAHRA; ARDIANSYAH, 2026) & 2026 & Medição & Latência e \textit{throughput} em APIs REST (H/3 vs H/2) \\
\hline
\end{tabular}
\fonte{Fonte: elaborado pelo autor. Tipo indica natureza predominante do documento (medição experimental, revisão ou análise comparativa).}
\end{table}

\subsection{Desempenho e latência}

No eixo de desempenho, predominam estudos de medição (TRIPATHI; MIRAZ; JOSHI, 2023; ZAHRA; ARDIANSYAH, 2026; YU; BENSON, 2021; PERNA et al., 2022) e revisões que sintetizam limitações (ZHANG et al., 2025). Tripathi, Miraz e Joshi (2023) compararam latência de APIs RESTful em HTTP/3 e HTTP/1.1, observando ganhos de eficiência de comunicação. Zhang et al. (2025) reúnem mecanismos de melhoria e restrições de desempenho. Zahra e Ardiansyah (2026) evidenciam a dependência de contexto: sob rede severamente degradada, relatam redução de latência de até 31,2\% com HTTP/3 frente ao HTTP/2, em especial em operações GET com alta concorrência, enquanto em condições ideais, com alta concorrência e payload médio, o HTTP/2 pode superar o HTTP/3 --- resultado que orienta a matriz experimental com netem e variação de payload neste TCC.

\subsection{Adoção e implementação}

Estudos de medição em larga escala (PERNA et al., 2022) e a revisão de Koch, Falowo e Elrod (2024) indicam adoção crescente do HTTP/3, com ganhos práticos que variam conforme as condições de rede e a infraestrutura (PERNA et al., 2022). Koch, Falowo e Elrod (2024) discutem barreiras de implementação e adoção do HTTP/3, em especial o suporte incompleto ao tráfego UDP em aplicações e serviços, ponto que motiva a escolha de proxy reverso no artefato Node.js.

\subsection{Segurança}

Chatzoglou et al. (2023) conduzem análise comparativa de segurança entre HTTP/3 e HTTP/2, com implicações para APIs expostas publicamente, terminação TLS na borda e configuração de certificados em proxies - aspectos relevantes ao NGINX utilizado no estudo de caso.

\subsection{Migração e arquitetura}

Yu e Benson (2021) mostram, por medição em produção, que o desempenho depende do algoritmo de congestionamento e da implementação do servidor - fator relevante também em arquiteturas com terminação na borda. Muñoz et al. (2025) comparam HTTP/3 e HTTP/2 com integração de proxy, abordagem alinhada à estratégia adotada neste TCC. As normas (QUIC Working Group - IETF, 2021; IETF, 2022) fundamentam semanticamente essa migração, mas não substituem evidência empírica sobre Node.js.

\section{Conclusão}

A revisão sistemática sintetizou nove estudos empíricos ou de revisão (Tabela 6), complementados pela base normativa das RFCs (Tabela 5), sobre desempenho, adoção, segurança e migração de HTTP/3 e QUIC em serviços web. Na Tabela 6, apenas Tripathi, Miraz e Joshi (2023) e Zahra e Ardiansyah (2026) tratam explicitamente de APIs REST. Nenhum dos estudos incluídos investiga a integração de HTTP/3 com Node.js por meio de proxy reverso em cenário controlado.

Essa lacuna motiva o estudo de caso experimental dos capítulos seguintes. O Capítulo 4 descreve a implementação do ambiente de testes, enquanto o Capítulo 5 apresenta os resultados obtidos sobre esse artefato e o Capítulo 6 sintetiza os achados e suas implicações para a adoção de HTTP/3 em APIs Node.js.

